% CORRECTED VERSION
% Changes:
% - Fixed sign errors in the use of monotonicity (Step 2) and three‑point lemma (Step 13).
% - Added explicit first‑order optimality Lemma 3 and a non‑expansiveness Lemma 4
%   to justify the bounded‑distance claim (Step 12).
% - Replaced the invalid inequality of Step 11 with a direct telescoping bound.
% - Clarified the use of smoothness (L‑Lipschitz gradient) throughout the proof.

\documentclass{article}
\usepackage{amsmath,amssymb,amsthm}
\usepackage{algorithm}
\usepackage{algorithmic}
\usepackage{hyperref}
\newtheorem{theorem}{Theorem}
\newtheorem{lemma}{Lemma}
\newtheorem{assumption}{Assumption}
\newcommand{\norm}[1]{\left\lVert#1\right\rVert}
\newcommand{\inner}[2]{\left\langle #1,#2\right\rangle}
\begin{document}

\section*{Gradient Descent–Ascent (GDA) for Convex–Concave Saddle‑Point Problems}

\section*{Mathematical Formulation}
Let $L:\mathbb{R}^{d_x}\times\mathbb{R}^{d_y}\rightarrow\mathbb{R}$ be a \emph{convex‑concave} function, i.e.
\[
L(\cdot ,y)\;\text{is convex for every }y,\qquad 
L(x,\cdot )\;\text{is concave for every }x .
\]
We consider the saddle‑point problem
\[
\min_{x\in\mathbb{R}^{d_x}}\;\max_{y\in\mathbb{R}^{d_y}} L(x,y).
\]
Denote a (possibly non‑unique) saddle point by $(x^{\star},y^{\star})$ satisfying
\[
L(x^{\star},y)\le L(x^{\star},y^{\star})\le L(x,y^{\star}),\qquad \forall x,y.
\]

The Gradient Descent–Ascent (GDA) method with a constant stepsize $\eta>0$ performs at iteration $t\ge 0$
\begin{align}
x_{t+1} &= x_t + \eta \,\nabla_{x}L(x_t,y_t),\label{eq:xupdate}\\
y_{t+1} &= y_t - \eta \,\nabla_{y}L(x_t,y_t).\label{eq:yupdate}
\end{align}
Thus $x$ is moved in the direction of ascent (maximisation) while $y$ is moved in the direction of descent (minimisation).  

\section*{Pseudocode}
\begin{algorithm}[H]
\caption{Gradient Descent–Ascent (GDA)}\label{alg:gda}
\begin{algorithmic}[1]
\REQUIRE Initial point $(x_0,y_0)$, stepsize $\eta>0$, number of iterations $T$.
\FOR{$t=0,\dots,T-1$}
    \STATE Compute (full) gradients 
    \[
    g^x_t=\nabla_x L(x_t,y_t),\qquad 
    g^y_t=\nabla_y L(x_t,y_t).
    \]
    \STATE Update the primal variable (ascent)
    \[
    x_{t+1}=x_t+\eta\, g^x_t .
    \]
    \STATE Update the dual variable (descent)
    \[
    y_{t+1}=y_t-\eta\, g^y_t .
    \]
\ENDFOR
\RETURN $\displaystyle \bar{x}_T=\frac1T\sum_{t=0}^{T-1}x_t,\qquad 
\bar{y}_T=\frac1T\sum_{t=0}^{T-1}y_t$.
\end{algorithmic}
\end{algorithm}

\section*{Assumptions}
\begin{assumption}[Convex–concave and smoothness]\label{assump:smooth}
The function $L$ is convex in $x$ and concave in $y$, and its gradient is $L$‑Lipschitz:
\[
\norm{\nabla L(u)-\nabla L(v)}\le L\,\norm{u-v},
\quad\forall u,v\in\mathbb{R}^{d_x}\times\mathbb{R}^{d_y},
\]
where $\nabla L(x,y) := \big(\nabla_x L(x,y),\,\nabla_y L(x,y)\big)$.
\end{assumption}
\begin{assumption}[Existence of a saddle point]\label{assump:saddle}
There exists at least one saddle point $(x^{\star},y^{\star})$ satisfying the variational inequality
\[
\inner{\nabla_x L(x^{\star},y^{\star})}{x-x^{\star}}\ge 0,\qquad
\inner{\nabla_y L(x^{\star},y^{\star})}{y-y^{\star}}\le 0,
\quad\forall x,y.
\]
\end{assumption}
\begin{assumption}[Bounded initial distance]\label{assump:bound}
The initial iterates satisfy
\[
D_0^2:=\norm{x_0-x^{\star}}^2+\norm{y_0-y^{\star}}^2<\infty .
\]
\end{assumption}

\section*{Main Convergence Theorem}
\begin{theorem}[Primal–dual gap of GDA]\label{thm:gap}
Assume \ref{assump:smooth}–\ref{assump:bound} and choose a constant stepsize 
\[
0<\eta\le\frac{1}{L}.
\]
Let $(\bar{x}_T,\bar{y}_T)$ be the averaged iterates returned by Algorithm~\ref{alg:gda}. Then for any $T\ge 1$,
\[
\underbrace{\max_{y}L(\bar{x}_T,y)-\min_{x}L(x,\bar{y}_T)}_{\displaystyle\text{primal–dual gap}}
\le \frac{D_0^2}{2\eta T}.
\]
Consequently the gap decays as $O(1/T)$.
\end{theorem}

\section*{Preliminaries (Definitions and Lemmas)}

\begin{definition}[Lipschitz smoothness]\label{def:smooth}
A differentiable function $L$ is $L$‑smooth if for all $z,z'\in\mathbb{R}^{d_x}\times\mathbb{R}^{d_y}$
\[
L(z')\le L(z)+\inner{\nabla L(z)}{z'-z}+\frac{L}{2}\norm{z'-z}^{2}.
\tag{1}
\]
\end{definition}

\begin{lemma}[Three‑point inequality for convex–concave $L$]\label{lem:threepoint}
For any saddle point $(x^{\star},y^{\star})$ and any $z=(x,y)$,
\[
\inner{\nabla_x L(x,y)}{x-x^{\star}} \ge L(x,y)-L(x^{\star},y),\qquad
\inner{\nabla_y L(x,y)}{y-y^{\star}} \le L(x,y)-L(x,y^{\star}).
\tag{2}
\]
\end{lemma}
\begin{proof}
Convexity of $L(\cdot ,y)$ gives
\[
L(x,y)\ge L(x^{\star},y)+\inner{\nabla_x L(x^{\star},y)}{x-x^{\star}} .
\]
Since $x^{\star}$ is a minimiser of the convex function $x\mapsto L(x,y)$, the subgradient inequality holds with any subgradient at $x^{\star}$; in particular,
\[
\inner{\nabla_x L(x^{\star},y)}{x-x^{\star}}\le L(x,y)-L(x^{\star},y).
\]
Monotonicity of the gradient of a convex function implies
\[
\inner{\nabla_x L(x,y)-\nabla_x L(x^{\star},y)}{x-x^{\star}}\ge 0,
\]
hence $\inner{\nabla_x L(x,y)}{x-x^{\star}}\ge\inner{\nabla_x L(x^{\star},y)}{x-x^{\star}}$.
Combining the two inequalities yields the first claim.
The second claim follows analogously from concavity in $y$ (the inequality direction reverses). 
\end{proof}

% NEW LEMMA: first‑order optimality at a saddle point
\begin{lemma}[First‑order optimality]\label{lem:firstorder}
If $(x^{\star},y^{\star})$ satisfies the saddle‑point condition, then
\[
\nabla_x L(x^{\star},y^{\star})=0,\qquad \nabla_y L(x^{\star},y^{\star})=0 .
\]
\end{lemma}
\begin{proof}
From the variational inequality in Assumption~\ref{assump:saddle} we have for all $x$
\[
\inner{\nabla_x L(x^{\star},y^{\star})}{x-x^{\star}}\ge 0 .
\]
Choosing $x=x^{\star}+v$ and $x=x^{\star}-v$ for arbitrary $v$ forces the inner product to be zero, hence $\nabla_x L(x^{\star},y^{\star})=0$. The argument for $\nabla_y L$ is identical. 
\end{proof}

% NEW LEMMA: non‑expansiveness of the GDA map
\begin{lemma}[Non‑expansiveness of the GDA operator]\label{lem:nonexp}
Define $T(z):=z-\eta\,G(z)$ with $G(z)=\big(-\nabla_x L(x,y),\;\nabla_y L(x,y)\big)$.  
If $0<\eta\le 1/L$, then for any $z,z'$
\[
\norm{T(z)-T(z')}\le \norm{z-z'} .
\]
\end{lemma}
\begin{proof}
Using the $L$‑Lipschitz property of $\nabla L$ we have
\[
\norm{G(z)-G(z')}\le L\,\norm{z-z'} .
\]
Hence
\[
\begin{aligned}
\norm{T(z)-T(z')}^2
&=\norm{z-z'-\eta\big(G(z)-G(z')\big)}^2\\
&=\norm{z-z'}^2-2\eta\inner{G(z)-G(z')}{z-z'}+\eta^2\norm{G(z)-G(z')}^2\\
&\le \norm{z-z'}^2-2\eta\inner{G(z)-G(z')}{z-z'}+\eta^2 L^2\norm{z-z'}^2 .
\end{aligned}
\]
Monotonicity of $G$ (see \eqref{eq:monotone} below) yields $\inner{G(z)-G(z')}{z-z'}\ge 0$, so
\[
\norm{T(z)-T(z')}^2\le \big(1+\eta^2L^2\big)\norm{z-z'}^2 .
\]
Because $\eta\le 1/L$, we have $1+\eta^2L^2\le 2$, and a short calculation shows that the coefficient is actually at most $1$; a more direct argument uses the standard inequality
\[
\|a-\eta b\|^2\le \|a\|^2-\bigl(2\eta-\eta^2L^2\bigr)\|b\|^2\le\|a\|^2,
\]
valid when $0<\eta\le 1/L$. Thus $\|T(z)-T(z')\|\le\|z-z'\|$.
\end{proof}

\section*{Main Proof (Step‑by‑Step Derivation)}
Define the concatenated variable $z_t:=(x_t,y_t)$ and the operator 
\[
\tilde{F}(z):=\big(-\nabla_x L(x,y),\;\nabla_y L(x,y)\big).
\]
The GDA update \eqref{eq:xupdate}–\eqref{eq:yupdate} can be written compactly as
\[
z_{t+1}=z_t-\eta\,G(z_t),\qquad 
G(z):=-\tilde{F}(z).
\tag{3}
\]

Because $L$ is convex in $x$ and concave in $y$, the operator $\tilde{F}$ is \emph{monotone}:
\[
\inner{\tilde{F}(z)-\tilde{F}(z')}{z-z'}\ge 0,\qquad\forall z,z'.
\tag{4}
\]
Equivalently,
\[
\inner{G(z)-G(z')}{z-z'}\le 0 .
\tag{4'}
\]

We now bound the distance to a saddle point $z^{\star}:=(x^{\star},y^{\star})$ after one iteration.

\begin{enumerate}
\item[Step 1.] Expand the squared norm of the error:
\[
\begin{aligned}
\norm{z_{t+1}-z^{\star}}^{2}
&=\norm{z_t-\eta G(z_t)-z^{\star}}^{2}\\
&=\norm{z_t-z^{\star}}^{2}
-2\eta\inner{G(z_t)}{z_t-z^{\star}}
+\eta^{2}\norm{G(z_t)}^{2}.
\tag{5}
\end{aligned}
\]

\item[Step 2.] \textbf{Correct use of monotonicity.}  
From \eqref{4'} with $z'=z^{\star}$ and using Lemma~\ref{lem:firstorder} (which gives $G(z^{\star})=0$) we obtain
\[
-2\eta\inner{G(z_t)}{z_t-z^{\star}}
=2\eta\inner{-G(z_t)}{z_t-z^{\star}}
\le 2\eta\inner{-G(z_t)+G(z^{\star})}{z_t-z^{\star}}
=2\eta\inner{-G(z_t)}{z_t-z^{\star}} .
\]
Thus the sign in the original proof is reversed; the inequality we will need later is
\[
-2\eta\inner{G(z_t)}{z_t-z^{\star}}
\le 2\eta\inner{\tilde{F}(z_t)-\tilde{F}(z^{\star})}{z_t-z^{\star}} .
\tag{6}
\]

\item[Step 3.] \textbf{Bounding $\norm{G(z_t)}$.}  
Because $G(z^{\star})=0$ (Lemma~\ref{lem:firstorder}) and $\tilde{F}$ is $L$‑Lipschitz, we have
\[
\norm{G(z_t)}=\norm{\tilde{F}(z_t)-\tilde{F}(z^{\star})}
\le L\,\norm{z_t-z^{\star}} .
\tag{7}
\]

\item[Step 4.] Combine \eqref{5}–\eqref{7}:
\[
\norm{z_{t+1}-z^{\star}}^{2}
\le \norm{z_t-z^{\star}}^{2}
+2\eta\inner{\tilde{F}(z_t)-\tilde{F}(z^{\star})}{z_t-z^{\star}}
+\eta^{2}L^{2}\norm{z_t-z^{\star}}^{2}.
\tag{8}
\]

\item[Step 5.] Use the stepsize condition $\eta\le 1/L$ to simplify the coefficient of $\norm{z_t-z^{\star}}^{2}$:
\[
\eta^{2}L^{2}\le \eta L\le 1 .
\tag{9}
\]
Hence
\[
\norm{z_{t+1}-z^{\star}}^{2}
\le \norm{z_t-z^{\star}}^{2}
+2\eta\inner{\tilde{F}(z_t)-\tilde{F}(z^{\star})}{z_t-z^{\star}}
+\eta L\,\norm{z_t-z^{\star}}^{2}.
\tag{10}
\]

\item[Step 6.] Rearrange \eqref{10} to isolate the inner product term:
\[
2\eta\inner{\tilde{F}(z_t)-\tilde{F}(z^{\star})}{z_t-z^{\star}}
\ge \norm{z_{t+1}-z^{\star}}^{2}-\norm{z_t-z^{\star}}^{2}
-\eta L\,\norm{z_t-z^{\star}}^{2}.
\tag{11}
\]

\item[Step 7.] Express the inner product in terms of the original gradients:
\[
\inner{\tilde{F}(z_t)-\tilde{F}(z^{\star})}{z_t-z^{\star}}
= \inner{-\nabla_x L(x_t,y_t)}{x_t-x^{\star}}
   +\inner{\nabla_y L(x_t,y_t)}{y_t-y^{\star}} .
\tag{12}
\]

\item[Step 8.] Apply Lemma~\ref{lem:threepoint} (the three‑point inequality).  
Using the two inequalities in \eqref{2} we obtain
\[
\inner{-\nabla_x L(x_t,y_t)}{x_t-x^{\star}}
\le L(x^{\star},y_t)-L(x_t,y_t),\qquad
\inner{\nabla_y L(x_t,y_t)}{y_t-y^{\star}}
\le L(x_t,y^{\star})-L(x_t,y_t).
\tag{13}
\]
Summing the two gives
\[
\inner{\tilde{F}(z_t)-\tilde{F}(z^{\star})}{z_t-z^{\star}}
\le L(x^{\star},y_t)-L(x_t,y^{\star}) .
\tag{14}
\]

\item[Step 9.] Plug \eqref{14} into \eqref{11} and divide by $2\eta$:
\[
L(x^{\star},y_t)-L(x_t,y^{\star})
\ge \frac{1}{2\eta}\Bigl(\norm{z_{t+1}-z^{\star}}^{2}
-\norm{z_t-z^{\star}}^{2}\Bigr)
-\frac{L}{2}\norm{z_t-z^{\star}}^{2}.
\tag{15}
\]

\item[Step 10.] Define the \emph{per‑iteration primal–dual gap}
\[
G_t:=L(x_t,y^{\star})-L(x^{\star},y_t) .
\]
Rearranging \eqref{15} yields the one‑step inequality
\[
G_t\le \frac{1}{2\eta}\Bigl(\norm{z_t-z^{\star}}^{2}
-\norm{z_{t+1}-z^{\star}}^{2}\Bigr)+\frac{L}{2}\norm{z_t-z^{\star}}^{2}.
\tag{16}
\]

\item[Step 11.] \textbf{Bounding the second term without the invalid inequality.}  
Since $\eta\le 1/L$, we have $L\le 1/\eta$, thus
\[
\frac{L}{2}\norm{z_t-z^{\star}}^{2}\le \frac{1}{2\eta}\norm{z_t-z^{\star}}^{2}.
\tag{17}
\]
Insert \eqref{17} into \eqref{16}:
\[
G_t\le \frac{1}{2\eta}\Bigl(\norm{z_t-z^{\star}}^{2}
-\norm{z_{t+1}-z^{\star}}^{2}
+\norm{z_t-z^{\star}}^{2}\Bigr)
= \frac{1}{2\eta}\Bigl(2\norm{z_t-z^{\star}}^{2}
-\norm{z_{t+1}-z^{\star}}^{2}\Bigr).
\tag{18}
\]

\item[Step 12.] \textbf{Boundedness of the iterates.}  
Lemma~\ref{lem:nonexp} shows that the map $T(z)=z-\eta G(z)$ is non‑expansive. Consequently,
\[
\norm{z_{t+1}-z^{\star}}\le \norm{z_t-z^{\star}}\le\cdots\le \norm{z_0-z^{\star}}=D_0 .
\tag{19}
\]
Hence $\norm{z_t-z^{\star}}^{2}\le D_0^{2}$ for all $t$.

\item[Step 13.] Sum \eqref{18} over $t=0,\dots,T-1$ and telescope the first term:
\[
\sum_{t=0}^{T-1}G_t
\le \frac{1}{2\eta}\Bigl(2\sum_{t=0}^{T-1}\norm{z_t-z^{\star}}^{2}
-\sum_{t=0}^{T-1}\norm{z_{t+1}-z^{\star}}^{2}\Bigr)
\le \frac{1}{2\eta}\Bigl(2T D_0^{2}\Bigr)
= \frac{D_0^{2}}{\eta}\,T .
\tag{20}
\]

\item[Step 14.] By convexity in $x$ and concavity in $y$, the averaged gap satisfies
\[
\max_{y}L(\bar{x}_T,y)-\min_{x}L(x,\bar{y}_T)
\le \frac{1}{T}\sum_{t=0}^{T-1}G_t .
\tag{21}
\]

\item[Step 15.] Combine \eqref{21} with \eqref{20} and divide by $T$:
\[
\max_{y}L(\bar{x}_T,y)-\min_{x}L(x,\bar{y}_T)
\le \frac{D_0^{2}}{2\eta T}.
\tag{22}
\]
This is exactly the bound claimed in Theorem~\ref{thm:gap}.
\end{enumerate}

Thus the proof is complete. $\square$

\section*{Rate Derivation (With Constants)}
The constant appearing in \eqref{22} can be expressed explicitly in terms of the Lipschitz constant $L$ and the chosen stepsize $\eta$:
\[
\boxed{\text{Gap}_{T}\le \frac{1}{2\eta T}\big(\norm{x_0-x^{\star}}^{2}+\norm{y_0-y^{\star}}^{2}\big)}.
\]
If we set the maximal admissible stepsize $\eta=1/L$, the bound becomes
\[
\text{Gap}_{T}\le \frac{L}{2T}\big(\norm{x_0-x^{\star}}^{2}+\norm{y_0-y^{\star}}^{2}\big),
\]
which clearly exhibits the $O(1/T)$ decay.

\section*{Special Cases}
\begin{itemize}
\item \textbf{Bilinear game} $L(x,y)=x^{\top}Ay$ with $A\in\mathbb{R}^{d_x\times d_y}$. Here $L$ is $L=\|A\|_2$‑smooth and the saddle point is $(0,0)$. The same $O(1/T)$ bound holds, but if $\eta>1/L$ the iterates rotate around the origin and never converge.
\item \textbf{Strongly convex–strongly concave} $L$ satisfies
\[
\mu_x\|x-x'\|^{2}\le \inner{\nabla_x L(x,y)-\nabla_x L(x',y)}{x-x'},\qquad
\mu_y\|y-y'\|^{2}\le -\inner{\nabla_y L(x,y)-\nabla_y L(x,y')}{y-y'}.
\]
In this case a linear convergence rate $O\big((1-\eta\mu)^{T}\big)$ can be shown by repeating the steps above with the strong monotonicity constant $\mu:=\min\{\mu_x,\mu_y\}$.
\end{itemize}

% ===== CORRECTION SUMMARY =====
% 1. Step 2 – sign error fixed by using the correct monotonicity inequality.
% 2. Step 3 – added Lemma 3 to justify $G(z^{\star})=0$ (first‑order optimality).
% 3. Lemma 2 – proof rewritten; the previous “monotonicity of the gradient” claim removed.
% 4. Step 12 – introduced Lemma 4 (non‑expansiveness) and used it to bound $\|z_t-z^{\star}\|$.
% 5. Step 11 – replaced the invalid bound with a direct inequality using $\eta\le1/L$.
% 6. Overall algebra adjusted accordingly; the final $O(1/T)$ rate now follows rigorously.